\begin{table}[H]
    \centering
    \caption{The three maneuvers of Fig.~\ref{fig:4dof_trajectory_riley}
    flown from $1000$ Latin-hypercube samples of the stalled-entry set
    $\gamma_0 \in [-30^\circ, 0]$, $V_0/V_s \in [0.40, 0.90]$,
    $\alpha_0 \in [14^\circ, 20^\circ]$, $q_0 \in [-20, 20]^\circ$/s,
    through the engine lag. The scripted pulls are capped to the pull
    authority the optimum commands at each entry, so the comparison
    isolates the procedure rather than the pilot's aggressiveness.
    Percentiles describe the uniform sampling box, not an operational
    entry distribution; the per-entry ranking reported below is
    distribution-free.}
    \label{tab:montecarlo_scripted}
    \setlength{\tabcolsep}{14pt}\renewcommand{\arraystretch}{1.15}
    \begin{tabular}{lrrrr}
        \hline
        & \multicolumn{2}{c}{$\Delta h$ (m)} & \multicolumn{2}{c}{excess over optimum (m)} \\
        Maneuver & median & 5th pct & median & 95th pct \\
        \hline
        DP optimum & $-61.73$ & $-91.72$ & --- & --- \\
        Scripted CAA (power with the push) & $-75.15$ & $-109.48$ & $13.47$ & $17.92$ \\
        Scripted FAA (power after unstall) & $-80.39$ & $-123.32$ & $18.18$ & $32.76$ \\
        \hline
    \end{tabular}
    \\[2pt]
    \footnotesize The CAA sequence loses less than the FAA sequence at $100.0\%$ of the sampled entries (median advantage $4.83$\,m, smallest $0.17$\,m); $999$ of $1000$ entries closed the recovery under every maneuver.
\end{table}
